\chapter{Evolution}
\label{chap:evolution}

This chapter introduces an evolutionary approach for improving the quality of packing solutions. 
Instead of designing a single fixed packing strategy, the algorithm evolves a population of candidate solutions and gradually improves them over time.

The main idea of evolutionary algorithms is to work with a set of candidate solutions, called \emph{individuals}, which are iteratively modified using operators inspired by natural evolution. 
In each generation, the algorithm evaluates the quality of all individuals, selects the more promising ones, and generates new candidate solutions through various evolutionary operators.
Through repeated iterations, the population is gradually refined and the algorithm searches for individuals that produce better packing solutions \cite{eiben2016}.

\section{Representation of an Individual}

An individual in the evolutionary algorithm represents the packing rules in their vector form, uniquely determining a complete packing solution, as defined in Section~\ref{sec:vector-packing-rules}.

\section{Fitness Evaluation}

The fitness of an individual corresponds directly to the fitness of the packing rules it encodes, as defined in Section~\ref{sec:fitness-packing-rules}.
Individuals representing better packing rules therefore obtain higher fitness values and are more likely to persist in the population in subsequent generations.

\section{Elitist Genetic Algorithm}

The evolutionary process used in this work is based on the elitist genetic algorithm introduced by Gonçalves and Resende \cite{goncalves2013}. 
The algorithm parameters are summarized in Table~\ref{tab:ega-hyperparameters}, and the full pseudocode is given in Figure~\ref{alg:ega-pseudocode}. 
In each generation, the population is divided into two groups: \emph{elite individuals} and \emph{non-elite individuals}. 
The elite group consists of the individuals with the highest fitness values, and its size is determined by a predefined parameter.

\subsection{Selection of Parents}

After dividing the population into elite and non-elite individuals, the same number of individuals is randomly selected from both groups. 
These selected individuals form the parent population used for generating new candidate solutions.

\subsection{Uniform Crossover}

New individuals are generated by applying a \emph{uniform crossover} to pairs of elite and non-elite parents. 
This type of crossover is called \emph{uniform} because each gene of the offspring is independently chosen from one of the two parents with a given probability, rather than swapping large contiguous segments of the genome. 
The probability that a gene is inherited from the elite parent determines how strongly the elite individuals influence the new generation.

\subsection{New Population Construction}

The new population in each generation consists of three parts:

\begin{itemize}
    \item the elite individuals, which are copied directly to the next generation,
    \item the offspring generated by crossover,
    \item a set of newly generated random individuals.
\end{itemize}

This strategy combines exploitation by preserving the best individuals and exploration by introducing new random individuals.
The relative sizes of these groups are determined by the parameters of the algorithm.

\subsection{Mutation}

In addition to the original algorithm, mutation is applied to the offspring generated by crossover. 
In this approach, each gene in the packing rules vector of an offspring individual is, with a certain probability $p_M$, replaced by a new random value drawn from its valid range, an interval $[0,1]$. 

The purpose of mutation is to test whether small random modifications can help the algorithm explore new regions of the solution space. 
In combination with newly generated random individuals, mutation may help discover better packing solutions that would not be reached by crossover alone.

\subsection{Memetic}

Another extension of the original algorithm is the memetic approach.
Here, a small number of \emph{hill-climbing} iterations is applied to the mutated individuals.

Hill climbing is a local search that applies small, mutation-like changes with a given probability.
For each iteration, a gene is randomly modified and the change is kept only if it improves fitness. 
This process is repeated for a fixed number of iterations $I_{HC}$.

We will test whether refining offspring in this way can improve convergence and solution quality by increasing their chances of becoming new elites.

\begin{table}[htbp]
\centering
\caption{Elitist Genetic Algorithm Hyperparameters}
\label{tab:ega-hyperparameters}
\begin{tabular}{|l|l|}
\hline
$M$ & population size \\ \hline
$G$ & maximum number of generations \\ \hline
$p_E$ & percentage of elite individuals in the population \\ \hline
$p_C$ & probability that a gene is inherited from the elite parent during uniform crossover \\ \hline
$p_M$ & probability that a gene in a packing rules vector is mutated in offspring \\ \hline
$p_R$ & percentage of newly generated random individuals each generation \\ \hline
$I_{HC}$ & number of hill-climbing iterations applied to each elite individual \\ \hline
$p_{HC}$ & probability that a gene is mutated during each hill-climbing step \\ \hline
\end{tabular}
\end{table}

\begin{algorithm}[htbp]
\centering
\begin{minipage}{0.95\linewidth}
\begin{algorithmic}[1]
\State \textbf{Input:} Parameters $M$, $G$, $p_E$, $p_R$, $p_C$, $p_M$, $I_{HC}$, $p_{HC}$
\State \textbf{Output:} Final population $\mathcal{P}$

\State $\mathcal{P} \gets \textbf{RandomPopulation}(M)$
\State Evaluate($\mathcal{P}$)

\State $N_{Elite} \gets \lceil p_E \cdot M \rceil$ \Comment{Number of elite individuals}
\State $N_{Random} \gets \lceil p_R \cdot M \rceil$ \Comment{Number of random individuals}
\State $N_{Selected} \gets M - N_{Elite} - N_{Random}$ \Comment{Number of parents selected for crossover}

\For{$g = 1$ to $G$} 

    \State Elite $\gets$ \textbf{ExtractElite}($\mathcal{P}$, $N_{Elite}$) \Comment{Select elite individuals}
    \State NonElite $\gets \mathcal{P} \setminus$ Elite \Comment{Remaining individuals}

    \State EliteParents $\gets$ \textbf{SelectRandom}(Elite, $N_{Selected}$) \Comment{Elite parents}
    \State NonEliteParents $\gets$ \textbf{SelectRandom}(NonElite, $N_{Selected}$) \Comment{Non-elite parents}

    \State Offspring $\gets$ \textbf{Crossover}(EliteParents, NonEliteParents, $p_C$) \Comment{Generate offspring}
    \State Mutated $\gets$ \textbf{Mutation}(Offspring, $p_M$) \Comment{Apply mutation}
    \State Improved $\gets$ \textbf{Memetic}(Mutated, $I_{HC}$, $p_{HC}$) \Comment{Local search improvement}

    \State RandomIndividuals $\gets$ \textbf{RandomPopulation}($N_{Random}$) \Comment{New random individuals}

    \State $\mathcal{P}' \gets$ Elite $\cup$ Improved $\cup$ RandomIndividuals

    \State Evaluate($\mathcal{P}'$)
    \State $\mathcal{P} \gets \mathcal{P}'$

\EndFor

\State \Return $\mathcal{P}$
\end{algorithmic}
\end{minipage}
\caption{Elitist Genetic Algorithm Pseudocode}
\label{alg:ega-pseudocode}
\end{algorithm}

\section{Probabilistic Hill Climbing}

As a simple baseline for comparison with the elitist genetic algorithm, we introduce probabilistic hill climbing. 
This method can be viewed as a simplified variant of simulated annealing \cite{eiben2016}, but instead of using a temperature parameter, it works with a probability of accepting a worse solution. 
This acceptance probability is gradually reduced over time by a geometric decay coefficient, but it is bounded below by a minimum acceptance threshold.

For simplicity, we reuse the terminology introduced for the evolutionary algorithm. 
At each step, a current individual is modified by a small random mutation, producing a neighboring candidate solution. 
If the candidate has a better fitness value, it is always accepted. 
If the candidate is worse, it may still be accepted with a certain probability $p_A$. 
After each generation, this probability is updated as
\[
p_A^{(g+1)} = \max(p_{\min}, \alpha \, p_A^{(g)}),
\]
where $\alpha \in (0,1)$ is the decay coefficient and $p_{\min}$ is the minimum acceptance threshold. 
Thus, the acceptance probability decreases geometrically over time, but never falls below the prescribed lower bound.

The main idea behind this approach is to balance exploration and exploitation. 
At the beginning of the search, a relatively high acceptance probability allows the algorithm to escape poor local optima and explore a wider portion of the solution space. 
Later, as the acceptance probability decreases, the algorithm becomes more selective and focuses on refining promising solutions. 
This behaviour often works well in practice because early exploration can reveal regions of the search space that would not be reached by purely greedy hill climbing, while later exploitation helps stabilize the search around high-quality solutions.

In the simplest setting, the algorithm repeatedly applies mutation and probabilistic acceptance to a population of individuals for a fixed number of generations. 
Although the method is conceptually close to local search, using multiple individuals in parallel makes it compatible with the same experimental framework as the evolutionary algorithm.

It is worth noting that if the acceptance probability is set to $p_A = 0$, the method reduces to classical hill climbing, where only strictly improving solutions are accepted. 
This variant will also be included in the experiments as a baseline for comparison.

The parameters of the probabilistic hill climbing baseline and its pseudocode are summarized in Table~\ref{tab:phc-hyperparameters} and Algorithm~\ref{alg:phc-pseudocode}.

\begin{table}[htbp]
\centering
\caption{Probabilistic Hill Climbing Hyperparameters}
\label{tab:phc-hyperparameters}
\begin{tabular}{|l|l|}
\hline
$M$ & population size \\ \hline
$G$ & maximum number of generations \\ \hline
$p_M$ & probability that a gene in an individual is mutated \\ \hline
$p_A^{(0)}$ & initial probability of accepting a worse solution \\ \hline
$\alpha$ & geometric decay coefficient for $p_A$ \\ \hline
$p_{\min}$ & minimum acceptance probability \\ \hline
\end{tabular}
\end{table}


\begin{algorithm}[htbp]
\centering
\begin{minipage}{0.95\linewidth}
\begin{algorithmic}[1]
\State \textbf{Input:} Parameters $M$, $G$, $p_M$, $p_A^{(0)}$, $\alpha$, $p_{\min}$
\State \textbf{Output:} Best individual found during the search

\State $\mathcal{P} \gets \textbf{RandomPopulation}(M)$
\State Evaluate($\mathcal{P}$)
\State $I^\star \gets \textbf{BestIndividual}(\mathcal{P})$
\State $p_A \gets p_A^{(0)}$

\For{$g = 1$ to $G$}
    \State $\mathcal{P}' \gets []$ \Comment{Initialize new population}

    \For{each individual $I \in \mathcal{P}$}
        \State $J \gets \textbf{Mutation}(I, p_M)$ \Comment{Generate a neighboring solution}

        \If{$F(J) \geq F(I)$}
            \State $I_{\text{new}} \gets J$ \Comment{Always accept improvement}
        \ElsIf{$\textbf{rand}() < p_A$}
            \State $I_{\text{new}} \gets J$ \Comment{Accept worse solution with probability $p_A$}
        \Else
            \State $I_{\text{new}} \gets I$ \Comment{Reject and keep current solution}
        \EndIf

        \State Add $I_{\text{new}}$ to $\mathcal{P}'$
    \EndFor

    \State Evaluate($\mathcal{P}'$) \Comment{Evaluate new population}
    \State $I^\star \gets \textbf{BestIndividual}(\mathcal{P}' \cup \{I^\star\})$ \Comment{Update global best solution}
    \State $\mathcal{P} \gets \mathcal{P}'$ \Comment{Move to next generation}
    \State $p_A \gets \max(p_{\min}, \alpha \, p_A)$ \Comment{Decrease acceptance probability with lower bound}
\EndFor

\State \Return $I^\star$
\end{algorithmic}
\end{minipage}
\caption{Probabilistic Hill Climbing Pseudocode}
\label{alg:phc-pseudocode}
\end{algorithm}